% conclusiones.tex — Conclusiones del documento.
% main.tex la incluye después de P6 y antes de Referencias.

\section*{Conclusiones}
\addcontentsline{toc}{section}{Conclusiones}

Las seis preguntas de este seminario abordan problemas distintos
—acoplamiento, elasticidad, resiliencia, granularidad, calidad y
reutilización—, pero llegan a la misma idea: ningún estilo domina a los
demás, porque cada decisión de diseño cambia una característica por otra. La
coreografía gana desacoplamiento y pierde visibilidad del flujo (P1);
Space-Based gana elasticidad y paga consistencia eventual (P2); el
\emph{Workflow Event Pattern} preserva la capacidad de respuesta del sistema
a costa de procesar los mensajes fuera de orden (P3); fijar la granularidad
por transacción de negocio reduce el acoplamiento a cambio de servicios más
grandes (P4); Service-Based Architecture ofrece consistencia más simple a
costa de agilidad de despliegue (P5); y los microservicios duplican código de
negocio para proteger su independencia de despliegue (P6).

Como resume la Primera Ley de la arquitectura de software, ``everything in
software architecture is a trade-off'' (Richards \& Ford, 2020, cap.~1): la
arquitectura correcta es la que elige conscientemente qué característica
sacrificar según las prioridades reales del sistema.
